\documentclass[12pt,a4paper]{article}
\usepackage[utf8]{vietnam}
\usepackage{amsmath,amssymb}
\usepackage{enumerate}
\usepackage{tabularx}
\usepackage{tikz}
\usepackage{xcolor}
\usepackage[top=0.45cm, bottom=0.45cm, left=0.8cm, right=0.8cm]{geometry}
\DeclareUnicodeCharacter{00B2}{\ensuremath{^2}}
\DeclareUnicodeCharacter{00B3}{\ensuremath{^3}}

\setlength{\parskip}{0pt}
\setlength{\parindent}{0pt}
\setlength{\fboxsep}{3.5pt}
\setlength{\fboxrule}{0.5pt}

% Dòng kẻ chấm mờ cự ly chuẩn viết tay tiểu học (18pt)
\newcommand{\dongke}[1]{%
  \begingroup%
  \setlength{\parskip}{0pt}%
  \baselineskip=18pt%
  \vspace{0.4mm}%
  \foreach \i in {1,...,#1}{%
    \noindent{\color{gray!55}\dotfill}\par%
  }%
  \vspace{0.2mm}%
  \endgroup%
}

% Khung bao quanh bài tập tự luận
\newcommand{\khungBaiTap}[1]{%
  \vspace{0.8mm}%
  \noindent\fbox{%
    \parbox{\dimexpr\linewidth-2\fboxsep-2\fboxrule\relax}{%
      \vspace{0.3mm}%
      #1%
      \vspace{0.3mm}%
    }%
  }%
  \vspace{0.5mm}%
}

% Ô bài tập đặt tính
\newcommand{\oDatTinh}[3][4]{%
  \vspace{0.3mm}%
  \noindent\textbf{#2} #3\par%
  \dongke{#1}%
}

% Câu trắc nghiệm dàn 4 đáp án trên 1 hàng ngang
\newcommand{\cauTN}[6]{%
  \noindent\textbf{Câu #1.} #2\par\vspace{0.1mm}%
  \noindent%
  \makebox[0.25\linewidth][l]{\textbf{A.} #3}%
  \makebox[0.25\linewidth][l]{\textbf{B.} #4}%
  \makebox[0.25\linewidth][l]{\textbf{C.} #5}%
  \makebox[0.25\linewidth][l]{\textbf{D.} #6}\par\vspace{0.6mm}%
}

% Câu trắc nghiệm dàn 2 đáp án / hàng
\newcommand{\cauTNhaiDong}[6]{%
  \noindent\textbf{Câu #1.} #2\par\vspace{0.1mm}%
  \noindent%
  \makebox[0.5\linewidth][l]{\textbf{A.} #3}%
  \makebox[0.5\linewidth][l]{\textbf{B.} #4}\par\vspace{0.1mm}%
  \noindent%
  \makebox[0.5\linewidth][l]{\textbf{C.} #5}%
  \makebox[0.5\linewidth][l]{\textbf{D.} #6}\par\vspace{0.6mm}%
}

\begin{document}
\pagestyle{empty}
\fontsize{10.2pt}{13.2pt}\selectfont

% =============================================================================
% ========================= MÃ ĐỀ: 401 ========================================
% =============================================================================

% -------------------- TRANG 1: MẶT TRƯỚC --------------------
\noindent
\begin{tabularx}{\linewidth}{|X|c|}
\hline
\textbf{PHIẾU BÀI TẬP CUỐI TUẦN -- TOÁN LỚP 4 (TUẦN 15)} & \textbf{MÃ ĐỀ: 401} \\
\textit{Trung tâm Toán Thầy Lĩnh -- Năm học: 2026 - 2027} & \textit{Thời gian: 45 phút} \\
\hline
\end{tabularx}

\vspace{1mm}

\noindent
\begin{tabularx}{\linewidth}{|X|p{3.5cm}|}
\hline
\rule[-2mm]{0pt}{7.5mm}\textbf{Họ và tên:} \dotfill & \textbf{Lớp:} \makebox[2cm]{\dotfill} \\
\hline
\end{tabularx}

\vspace{1mm}

\noindent
\begin{tabularx}{\linewidth}{|p{3cm}|X|}
\hline
\centering\textbf{ĐIỂM SỐ} & \centering\textbf{LỜI NHẬN XÉT CỦA THẦY CÔ} \tabularnewline
\hline
\rule{0pt}{1.1cm} & \tabularnewline
\hline
\end{tabularx}

\vspace{1mm}

\noindent\textbf{BẢNG ĐIỀN ĐÁP ÁN TRẮC NGHIỆM (Mỗi câu đúng 0,25 điểm):}
\vspace{0.6mm}

\begin{tabularx}{\linewidth}{|c|*{12}{>{\centering\arraybackslash}X|}}
\hline
\textbf{Câu} & \textbf{1} & \textbf{2} & \textbf{3} & \textbf{4} & \textbf{5} & \textbf{6} & \textbf{7} & \textbf{8} & \textbf{9} & \textbf{10} & \textbf{11} & \textbf{12} \\
\hline
\textbf{Đáp án} & & & & & & & & & & & & \\
\hline
\end{tabularx}

\vspace{1.2mm}

\noindent\textbf{\large PHẦN I. TRẮC NGHIỆM KHÁCH QUAN (3,0 điểm)}

\cauTN{1}{Hai đường thẳng song song là hai đường thẳng:}{Cắt nhau tại 1 điểm}{Cắt nhau và tạo thành góc vuông}{Không bao giờ cắt nhau}{Trùng khít lên nhau}
\cauTN{2}{Trong hình chữ nhật $ABCD$, cạnh song song với cạnh $AB$ là:}{Cạnh $BC$}{Cạnh $CD$}{Cạnh $AD$}{Cạnh $AC$}
\cauTN{3}{Hình chữ nhật có bao nhiêu cặp cạnh đối diện song song với nhau?}{$1$ cặp}{$2$ cặp}{$3$ cặp}{$4$ cặp}
\cauTN{4}{Hai đường thẳng cùng vuông góc với một đường thẳng thứ ba thì chúng:}{Vuông góc với nhau}{Song song với nhau}{Cắt nhau}{Trùng nhau}
\cauTN{5}{Hình vuông có các cặp cạnh đối diện:}{Cắt nhau}{Song song và bằng nhau}{Vuông góc và không bằng nhau}{Trùng nhau}
\cauTN{6}{Số thích hợp điền vào chỗ chấm $6\,\text{tấn}\,350\,\text{kg} = \dots\,\text{kg}$ là:}{$635$}{$6\,350$}{$60\,350$}{$63\,500$}
\cauTN{7}{Số thích hợp điền vào chỗ chấm $5\,\text{m}^2\,42\,\text{dm}^2 = \dots\,\text{dm}^2$ là:}{$542$}{$5\,042$}{$5\,420$}{$56$}
\cauTN{8}{Khoảng thời gian $2\,\text{phút}\,50\,\text{giây}$ bằng bao nhiêu giây?}{$150\,\text{giây}$}{$170\,\text{giây}$}{$180\,\text{giây}$}{$190\,\text{giây}$}
\cauTN{9}{Tìm số hạng chưa biết: $x + 45\,200 = 92\,650$:}{$x = 47\,450$}{$x = 47\,550$}{$x = 137\,850$}{$x = 46\,450$}
\cauTN{10}{Làm tròn số $563\,800$ đến hàng trăm nghìn ta được số:}{$500\,000$}{$600\,000$}{$560\,000$}{$570\,000$}
\cauTNhaiDong{11}{Tổng của hai số là $120$, hiệu của hai số là $30$. Số lớn là:}{$70$}{$75$}{$80$}{$85$}
\cauTNhaiDong{12}{Một mảnh đất hình chữ nhật có chu vi $160\,\text{m}$, chiều dài hơn chiều rộng $20\,\text{m}$. Chiều dài là:}{$40\,\text{m}$}{$50\,\text{m}$}{$60\,\text{m}$}{$30\,\text{m}$}

\vspace{0.4mm}
\noindent\textbf{\large PHẦN II. TỰ LUẬN (7,0 điểm)}

\noindent\textbf{Bài 1 (2,0 điểm). Đặt tính rồi tính:}
\vspace{0.4mm}

\noindent
\begin{tabularx}{\linewidth}{|X|X|X|X|}
\hline
\oDatTinh[4]{a)}{$458\,720 + 269\,548$} & 
\oDatTinh[4]{b)}{$845\,310 - 478\,625$} & 
\oDatTinh[4]{c)}{$16\,435 \times 5$} & 
\oDatTinh[4]{d)}{$58\,464 : 8$} \tabularnewline
\hline
\end{tabularx}

\newpage
% -------------------- TRANG 2: MẶT SAU --------------------
\noindent
\begin{tabularx}{\linewidth}{|X|c|}
\hline
\textbf{PHẦN TỰ LUẬN (TIẾP THEO) -- PHIẾU BÀI TẬP TUẦN 15} & \textbf{MÃ ĐỀ: 401} \\
\hline
\end{tabularx}

\vspace{0.8mm}

% BÀI 2: TÌM X (4 Ý)
\khungBaiTap{%
  \textbf{Bài 2 (1,5 điểm). Tìm $x$:}\par\vspace{0.6mm}
  \noindent
  \begin{minipage}[t]{0.485\linewidth}
    \textbf{a)} $x + 38\,490 = 85\,210$\par
    \dongke{3}
    \vspace{0.4mm}
    \textbf{c)} $94\,300 - x = 46\,580$\par
    \dongke{3}
  \end{minipage}\hfill
  \begin{minipage}[t]{0.485\linewidth}
    \textbf{b)} $x - 47\,620 = 76\,850$\par
    \dongke{3}
    \vspace{0.4mm}
    \textbf{d)} $x : 7 = 15\,248$\par
    \dongke{3}
  \end{minipage}
}

\vspace{0.8mm}

% BÀI 3: ĐỔI ĐƠN VỊ (4 Ý - 1 CỘT TỪ TRÊN XUỐNG)
\khungBaiTap{%
  \textbf{Bài 3 (1,0 điểm). Viết số thích hợp vào chỗ chấm:}\par\vspace{0.8mm}
  \noindent
  a) $7\,\text{tấn}\,450\,\text{kg} = \dotfill\,\text{kg}$\par\vspace{1.4mm}
  \noindent
  b) $8\,\text{m}^2\,36\,\text{dm}^2 = \dotfill\,\text{dm}^2$\par\vspace{1.4mm}
  \noindent
  c) $6\,\text{thế kỉ}\,25\,\text{năm} = \dotfill\,\text{năm}$\par\vspace{1.4mm}
  \noindent
  d) $4\,\text{phút}\,20\,\text{giây} = \dotfill\,\text{giây}$\par
}

\vspace{0.8mm}

% BÀI 4: BÀI TOÁN GIẢI (11 DÒNG KẺ)
\khungBaiTap{%
  \textbf{Bài 4 (1,5 điểm).} Một mảnh đất hình thang có đáy lớn $24$ m, đáy bé $16$ m, chiều cao $12$ m. Người ta dành $\frac{1}{4}$ diện tích mảnh đất để làm nhà, phần còn lại lát gạch. Biết tiền gạch và công lát là $180\,000$ đồng mỗi mét vuông. Tính số tiền lát gạch phần còn lại.
  \par\vspace{0.3mm}%
  \noindent\textit{Bài giải:}
  \dongke{11}
}

\vspace{0.8mm}

% BÀI 5: TOÁN NÂNG CAO (8 DÒNG KẺ)
\khungBaiTap{%
  \textbf{Bài 5 (1,0 điểm).} Tính bằng cách thuận tiện nhất:\par\vspace{0.6mm}\centerline{$Tính bằng cách thuận tiện: A = 0{,}1 + 0{,}2 + 0{,}3 + ... + 0{,}9$}
  \par\vspace{0.3mm}%
  \noindent\textit{Bài làm:}
  \dongke{8}
}

\newpage


% =============================================================================
% ========================= MÃ ĐỀ: 402 ========================================
% =============================================================================

% -------------------- TRANG 1: MẶT TRƯỚC --------------------
\noindent
\begin{tabularx}{\linewidth}{|X|c|}
\hline
\textbf{PHIẾU BÀI TẬP CUỐI TUẦN -- TOÁN LỚP 4 (TUẦN 15)} & \textbf{MÃ ĐỀ: 402} \\
\textit{Trung tâm Toán Thầy Lĩnh -- Năm học: 2026 - 2027} & \textit{Thời gian: 45 phút} \\
\hline
\end{tabularx}

\vspace{1mm}

\noindent
\begin{tabularx}{\linewidth}{|X|p{3.5cm}|}
\hline
\rule[-2mm]{0pt}{7.5mm}\textbf{Họ và tên:} \dotfill & \textbf{Lớp:} \makebox[2cm]{\dotfill} \\
\hline
\end{tabularx}

\vspace{1mm}

\noindent
\begin{tabularx}{\linewidth}{|p{3cm}|X|}
\hline
\centering\textbf{ĐIỂM SỐ} & \centering\textbf{LỜI NHẬN XÉT CỦA THẦY CÔ} \tabularnewline
\hline
\rule{0pt}{1.1cm} & \tabularnewline
\hline
\end{tabularx}

\vspace{1mm}

\noindent\textbf{BẢNG ĐIỀN ĐÁP ÁN TRẮC NGHIỆM (Mỗi câu đúng 0,25 điểm):}
\vspace{0.6mm}

\begin{tabularx}{\linewidth}{|c|*{12}{>{\centering\arraybackslash}X|}}
\hline
\textbf{Câu} & \textbf{1} & \textbf{2} & \textbf{3} & \textbf{4} & \textbf{5} & \textbf{6} & \textbf{7} & \textbf{8} & \textbf{9} & \textbf{10} & \textbf{11} & \textbf{12} \\
\hline
\textbf{Đáp án} & & & & & & & & & & & & \\
\hline
\end{tabularx}

\vspace{1.2mm}

\noindent\textbf{\large PHẦN I. TRẮC NGHIỆM KHÁCH QUAN (3,0 điểm)}

\cauTN{1}{Hai đường thẳng song song với nhau thì:}{Không bao giờ gặp nhau}{Cắt nhau tại trung điểm}{Vuông góc với nhau}{Tạo thành 4 góc vuông}
\cauTN{2}{Trong hình vuông $ABCD$, cạnh song song với cạnh $AD$ là:}{Cạnh $AB$}{Cạnh $BC$}{Cạnh $CD$}{Cạnh $AC$}
\cauTN{3}{Một hình chữ nhật có bao nhiêu cạnh?}{$2$ cạnh}{$3$ cạnh}{$4$ cạnh}{$6$ cạnh}
\cauTN{4}{Để vẽ đường thẳng đi qua điểm $M$ và song song với đường thẳng $a$, ta thường vẽ:}{Đường thẳng qua $M$ vuông góc với đường vuông góc với $a$}{Đường thẳng bất kì qua $M$}{Đường tròn qua $M$}{Đoạn thẳng nối $M$ với $a$}
\cauTN{5}{Cặp cạnh đối diện của hình chữ nhật có tính chất gì?}{Vuông góc với nhau}{Song song và bằng nhau}{Cắt nhau tại 1 điểm}{Không bằng nhau}
\cauTN{6}{Số thích hợp điền vào chỗ chấm $8\,\text{tấn}\,40\,\text{kg} = \dots\,\text{kg}$ là:}{$840$}{$8\,040$}{$8\,400$}{$84\,000$}
\cauTN{7}{Số thích hợp điền vào chỗ chấm $6\,\text{m}^2\,75\,\text{dm}^2 = \dots\,\text{dm}^2$ là:}{$675$}{$6\,075$}{$6\,750$}{$69$}
\cauTN{8}{Khoảng thời gian $3\,\text{phút}\,15\,\text{giây}$ bằng bao nhiêu giây?}{$185\,\text{giây}$}{$195\,\text{giây}$}{$205\,\text{giây}$}{$215\,\text{giây}$}
\cauTN{9}{Tìm $x$: $x - 36\,420 = 54\,890$:}{$x = 91\,310$}{$x = 18\,470$}{$x = 91\,210$}{$x = 90\,310$}
\cauTN{10}{Làm tròn số $674\,920$ đến hàng trăm nghìn ta được số:}{$600\,000$}{$700\,000$}{$670\,000$}{$680\,000$}
\cauTNhaiDong{11}{Tổng của hai số là $150$, hiệu của hai số là $50$. Số bé là:}{$45$}{$50$}{$55$}{$60$}
\cauTNhaiDong{12}{Một mảnh đất hình chữ nhật có chu vi $180\,\text{m}$, chiều dài hơn chiều rộng $22\,\text{m}$. Chiều rộng là:}{$34\,\text{m}$}{$56\,\text{m}$}{$35\,\text{m}$}{$40\,\text{m}$}

\vspace{0.4mm}
\noindent\textbf{\large PHẦN II. TỰ LUẬN (7,0 điểm)}

\noindent\textbf{Bài 1 (2,0 điểm). Đặt tính rồi tính:}
\vspace{0.4mm}

\noindent
\begin{tabularx}{\linewidth}{|X|X|X|X|}
\hline
\oDatTinh[4]{a)}{$534\,680 + 287\,495$} & 
\oDatTinh[4]{b)}{$912\,400 - 456\,835$} & 
\oDatTinh[4]{c)}{$17\,348 \times 4$} & 
\oDatTinh[4]{d)}{$65\,430 : 6$} \tabularnewline
\hline
\end{tabularx}

\newpage
% -------------------- TRANG 2: MẶT SAU --------------------
\noindent
\begin{tabularx}{\linewidth}{|X|c|}
\hline
\textbf{PHẦN TỰ LUẬN (TIẾP THEO) -- PHIẾU BÀI TẬP TUẦN 15} & \textbf{MÃ ĐỀ: 402} \\
\hline
\end{tabularx}

\vspace{0.8mm}

% BÀI 2: TÌM X (4 Ý)
\khungBaiTap{%
  \textbf{Bài 2 (1,5 điểm). Tìm $x$:}\par\vspace{0.6mm}
  \noindent
  \begin{minipage}[t]{0.485\linewidth}
    \textbf{a)} $x + 46\,380 = 92\,150$\par
    \dongke{3}
    \vspace{0.4mm}
    \textbf{c)} $96\,200 - x = 48\,730$\par
    \dongke{3}
  \end{minipage}\hfill
  \begin{minipage}[t]{0.485\linewidth}
    \textbf{b)} $x - 54\,370 = 82\,490$\par
    \dongke{3}
    \vspace{0.4mm}
    \textbf{d)} $x : 8 = 16\,342$\par
    \dongke{3}
  \end{minipage}
}

\vspace{0.8mm}

% BÀI 3: ĐỔI ĐƠN VỊ (4 Ý - 1 CỘT TỪ TRÊN XUỐNG)
\khungBaiTap{%
  \textbf{Bài 3 (1,0 điểm). Viết số thích hợp vào chỗ chấm:}\par\vspace{0.8mm}
  \noindent
  a) $6\,\text{tấn}\,850\,\text{kg} = \dotfill\,\text{kg}$\par\vspace{1.4mm}
  \noindent
  b) $9\,\text{m}^2\,48\,\text{dm}^2 = \dotfill\,\text{dm}^2$\par\vspace{1.4mm}
  \noindent
  c) $7\,\text{thế kỉ}\,15\,\text{năm} = \dotfill\,\text{năm}$\par\vspace{1.4mm}
  \noindent
  d) $5\,\text{phút}\,20\,\text{giây} = \dotfill\,\text{giây}$\par
}

\vspace{0.8mm}

% BÀI 4: BÀI TOÁN GIẢI (11 DÒNG KẺ)
\khungBaiTap{%
  \textbf{Bài 4 (1,5 điểm).} Một khu vườn hình chữ nhật có chiều dài $30$ m, chiều rộng $20$ m. Ở giữa vườn người ta đào một cái ao hình tam giác có đáy $8$ m và chiều cao $5$ m. Phần đất còn lại trồng cỏ, mỗi mét vuông cỏ hết $25\,000$ đồng. Tính số tiền trồng cỏ.
  \par\vspace{0.3mm}%
  \noindent\textit{Bài giải:}
  \dongke{11}
}

\vspace{0.8mm}

% BÀI 5: TOÁN NÂNG CAO (8 DÒNG KẺ)
\khungBaiTap{%
  \textbf{Bài 5 (1,0 điểm).} Tính bằng cách thuận tiện nhất:\par\vspace{0.6mm}\centerline{$Tính bằng cách thuận tiện: B = 1{,}25 \times 3{,}7 + 1{,}25 \times 6{,}3$}
  \par\vspace{0.3mm}%
  \noindent\textit{Bài làm:}
  \dongke{8}
}

\newpage


% =============================================================================
% ========================= MÃ ĐỀ: 403 ========================================
% =============================================================================

% -------------------- TRANG 1: MẶT TRƯỚC --------------------
\noindent
\begin{tabularx}{\linewidth}{|X|c|}
\hline
\textbf{PHIẾU BÀI TẬP CUỐI TUẦN -- TOÁN LỚP 4 (TUẦN 15)} & \textbf{MÃ ĐỀ: 403} \\
\textit{Trung tâm Toán Thầy Lĩnh -- Năm học: 2026 - 2027} & \textit{Thời gian: 45 phút} \\
\hline
\end{tabularx}

\vspace{1mm}

\noindent
\begin{tabularx}{\linewidth}{|X|p{3.5cm}|}
\hline
\rule[-2mm]{0pt}{7.5mm}\textbf{Họ và tên:} \dotfill & \textbf{Lớp:} \makebox[2cm]{\dotfill} \\
\hline
\end{tabularx}

\vspace{1mm}

\noindent
\begin{tabularx}{\linewidth}{|p{3cm}|X|}
\hline
\centering\textbf{ĐIỂM SỐ} & \centering\textbf{LỜI NHẬN XÉT CỦA THẦY CÔ} \tabularnewline
\hline
\rule{0pt}{1.1cm} & \tabularnewline
\hline
\end{tabularx}

\vspace{1mm}

\noindent\textbf{BẢNG ĐIỀN ĐÁP ÁN TRẮC NGHIỆM (Mỗi câu đúng 0,25 điểm):}
\vspace{0.6mm}

\begin{tabularx}{\linewidth}{|c|*{12}{>{\centering\arraybackslash}X|}}
\hline
\textbf{Câu} & \textbf{1} & \textbf{2} & \textbf{3} & \textbf{4} & \textbf{5} & \textbf{6} & \textbf{7} & \textbf{8} & \textbf{9} & \textbf{10} & \textbf{11} & \textbf{12} \\
\hline
\textbf{Đáp án} & & & & & & & & & & & & \\
\hline
\end{tabularx}

\vspace{1.2mm}

\noindent\textbf{\large PHẦN I. TRẮC NGHIỆM KHÁCH QUAN (3,0 điểm)}

\cauTN{1}{Hai đường thẳng song song với nhau thì có bao nhiêu điểm chung?}{$0$ điểm chung}{$1$ điểm chung}{$2$ điểm chung}{Vô số điểm chung}
\cauTN{2}{Trong hình chữ nhật $MNPQ$, cạnh nào song song với cạnh $PQ$?}{Cạnh $MN$}{Cạnh $NP$}{Cạnh $QM$}{Cạnh $MP$}
\cauTN{3}{Trong hình vuông, các cạnh đối diện có quan hệ:}{Vuông góc}{Song song}{Cắt nhau}{Trùng nhau}
\cauTN{4}{Đoạn thẳng nối hai điểm trên hai đường thẳng song song và vuông góc với chúng thì biểu thị:}{Góc nhọn}{Khoảng cách giữa hai đường thẳng}{Chu vi}{Diện tích}
\cauTN{5}{Hình thoi có các cặp cạnh đối diện:}{Vuông góc}{Song song}{Cắt nhau}{Không bằng nhau}
\cauTN{6}{Số thích hợp điền vào chỗ chấm $9\,\text{tấn}\,60\,\text{kg} = \dots\,\text{kg}$ là:}{$960$}{$9\,060$}{$9\,600$}{$96\,000$}
\cauTN{7}{Số thích hợp điền vào chỗ chấm $7\,\text{m}^2\,65\,\text{dm}^2 = \dots\,\text{dm}^2$ là:}{$765$}{$7\,065$}{$7\,650$}{$78$}
\cauTN{8}{Khoảng thời gian $4\,\text{phút}\,40\,\text{giây}$ bằng bao nhiêu giây?}{$240\,\text{giây}$}{$260\,\text{giây}$}{$280\,\text{giây}$}{$300\,\text{giây}$}
\cauTN{9}{Tìm $x$: $x + 53\,470 = 98\,210$:}{$x = 44\,740$}{$x = 44\,840$}{$x = 151\,680$}{$x = 43\,740$}
\cauTN{10}{Làm tròn số $785\,300$ đến hàng trăm nghìn ta được số:}{$700\,000$}{$800\,000$}{$780\,000$}{$790\,000$}
\cauTNhaiDong{11}{Tổng của hai số là $180$, hiệu của hai số là $40$. Số lớn là:}{$100$}{$110$}{$120$}{$130$}
\cauTNhaiDong{12}{Một mảnh vườn hình chữ nhật có chu vi $140\,\text{m}$, chiều dài hơn chiều rộng $18\,\text{m}$. Chiều dài là:}{$44\,\text{m}$}{$26\,\text{m}$}{$45\,\text{m}$}{$50\,\text{m}$}

\vspace{0.4mm}
\noindent\textbf{\large PHẦN II. TỰ LUẬN (7,0 điểm)}

\noindent\textbf{Bài 1 (2,0 điểm). Đặt tính rồi tính:}
\vspace{0.4mm}

\noindent
\begin{tabularx}{\linewidth}{|X|X|X|X|}
\hline
\oDatTinh[4]{a)}{$624\,750 + 298\,685$} & 
\oDatTinh[4]{b)}{$853\,210 - 478\,564$} & 
\oDatTinh[4]{c)}{$18\,245 \times 5$} & 
\oDatTinh[4]{d)}{$74\,656 : 8$} \tabularnewline
\hline
\end{tabularx}

\newpage
% -------------------- TRANG 2: MẶT SAU --------------------
\noindent
\begin{tabularx}{\linewidth}{|X|c|}
\hline
\textbf{PHẦN TỰ LUẬN (TIẾP THEO) -- PHIẾU BÀI TẬP TUẦN 15} & \textbf{MÃ ĐỀ: 403} \\
\hline
\end{tabularx}

\vspace{0.8mm}

% BÀI 2: TÌM X (4 Ý)
\khungBaiTap{%
  \textbf{Bài 2 (1,5 điểm). Tìm $x$:}\par\vspace{0.6mm}
  \noindent
  \begin{minipage}[t]{0.485\linewidth}
    \textbf{a)} $x + 48\,650 = 95\,340$\par
    \dongke{3}
    \vspace{0.4mm}
    \textbf{c)} $97\,400 - x = 52\,650$\par
    \dongke{3}
  \end{minipage}\hfill
  \begin{minipage}[t]{0.485\linewidth}
    \textbf{b)} $x - 56\,420 = 78\,910$\par
    \dongke{3}
    \vspace{0.4mm}
    \textbf{d)} $x : 9 = 14\,236$\par
    \dongke{3}
  \end{minipage}
}

\vspace{0.8mm}

% BÀI 3: ĐỔI ĐƠN VỊ (4 Ý - 1 CỘT TỪ TRÊN XUỐNG)
\khungBaiTap{%
  \textbf{Bài 3 (1,0 điểm). Viết số thích hợp vào chỗ chấm:}\par\vspace{0.8mm}
  \noindent
  a) $8\,\text{tấn}\,240\,\text{kg} = \dotfill\,\text{kg}$\par\vspace{1.4mm}
  \noindent
  b) $6\,\text{m}^2\,85\,\text{dm}^2 = \dotfill\,\text{dm}^2$\par\vspace{1.4mm}
  \noindent
  c) $8\,\text{thế kỉ}\,30\,\text{năm} = \dotfill\,\text{năm}$\par\vspace{1.4mm}
  \noindent
  d) $4\,\text{phút}\,55\,\text{giây} = \dotfill\,\text{giây}$\par
}

\vspace{0.8mm}

% BÀI 4: BÀI TOÁN GIẢI (11 DÒNG KẺ)
\khungBaiTap{%
  \textbf{Bài 4 (1,5 điểm).} Một mảnh đất gồm hình chữ nhật có chiều dài $18$ m, chiều rộng $12$ m và một hình tam giác ghép vào có đáy bằng chiều rộng hình chữ nhật, chiều cao $6$ m. Người ta lát gạch toàn bộ mảnh đất với giá $150\,000$ đồng mỗi mét vuông. Tính số tiền lát gạch.
  \par\vspace{0.3mm}%
  \noindent\textit{Bài giải:}
  \dongke{11}
}

\vspace{0.8mm}

% BÀI 5: TOÁN NÂNG CAO (8 DÒNG KẺ)
\khungBaiTap{%
  \textbf{Bài 5 (1,0 điểm).} Tính bằng cách thuận tiện nhất:\par\vspace{0.6mm}\centerline{$Tính bằng cách thuận tiện: C = 0{,}25 \times 0{,}4 \times 8 \times 12{,}5$}
  \par\vspace{0.3mm}%
  \noindent\textit{Bài làm:}
  \dongke{8}
}

\newpage


% =============================================================================
% ========================= MÃ ĐỀ: 404 ========================================
% =============================================================================

% -------------------- TRANG 1: MẶT TRƯỚC --------------------
\noindent
\begin{tabularx}{\linewidth}{|X|c|}
\hline
\textbf{PHIẾU BÀI TẬP CUỐI TUẦN -- TOÁN LỚP 4 (TUẦN 15)} & \textbf{MÃ ĐỀ: 404} \\
\textit{Trung tâm Toán Thầy Lĩnh -- Năm học: 2026 - 2027} & \textit{Thời gian: 45 phút} \\
\hline
\end{tabularx}

\vspace{1mm}

\noindent
\begin{tabularx}{\linewidth}{|X|p{3.5cm}|}
\hline
\rule[-2mm]{0pt}{7.5mm}\textbf{Họ và tên:} \dotfill & \textbf{Lớp:} \makebox[2cm]{\dotfill} \\
\hline
\end{tabularx}

\vspace{1mm}

\noindent
\begin{tabularx}{\linewidth}{|p{3cm}|X|}
\hline
\centering\textbf{ĐIỂM SỐ} & \centering\textbf{LỜI NHẬN XÉT CỦA THẦY CÔ} \tabularnewline
\hline
\rule{0pt}{1.1cm} & \tabularnewline
\hline
\end{tabularx}

\vspace{1mm}

\noindent\textbf{BẢNG ĐIỀN ĐÁP ÁN TRẮC NGHIỆM (Mỗi câu đúng 0,25 điểm):}
\vspace{0.6mm}

\begin{tabularx}{\linewidth}{|c|*{12}{>{\centering\arraybackslash}X|}}
\hline
\textbf{Câu} & \textbf{1} & \textbf{2} & \textbf{3} & \textbf{4} & \textbf{5} & \textbf{6} & \textbf{7} & \textbf{8} & \textbf{9} & \textbf{10} & \textbf{11} & \textbf{12} \\
\hline
\textbf{Đáp án} & & & & & & & & & & & & \\
\hline
\end{tabularx}

\vspace{1.2mm}

\noindent\textbf{\large PHẦN I. TRẮC NGHIỆM KHÁCH QUAN (3,0 điểm)}

\cauTN{1}{Khi kéo dài mãi về hai phía của hai đường thẳng song song, chúng sẽ:}{Cắt nhau ở điểm rất xa}{Không bao giờ cắt nhau}{Vuông góc với nhau}{Gặp nhau tại 2 điểm}
\cauTN{2}{Trong hình chữ nhật $ABCD$, cặp cạnh nào sau đây song song với nhau?}{$AB$ và $BC$}{$AD$ và $BC$}{$AB$ và $AD$}{$CD$ và $AD$}
\cauTN{3}{Đặc điểm nào đúng với hình chữ nhật?}{Có 4 góc vuông và 2 cặp cạnh đối diện song song}{Có 4 góc nhọn}{Có 4 cạnh bằng nhau}{Có 2 góc tù}
\cauTN{4}{Dùng ê-ke và thước kẻ để vẽ đường thẳng song song dựa trên quy tắc:}{Cùng tạo góc nhọn}{Cùng vuông góc với một đường thẳng}{Cùng đi qua một điểm}{Cùng cắt nhau}
\cauTN{5}{Một hình bình hành có bao nhiêu cặp cạnh đối diện song song?}{$1$ cặp}{$2$ cặp}{$3$ cặp}{$4$ cặp}
\cauTN{6}{Số thích hợp điền vào chỗ chấm $5\,\text{tấn}\,75\,\text{kg} = \dots\,\text{kg}$ là:}{$575$}{$5\,075$}{$5\,750$}{$57\,500$}
\cauTN{7}{Số thích hợp điền vào chỗ chấm $8\,\text{m}^2\,56\,\text{dm}^2 = \dots\,\text{dm}^2$ là:}{$856$}{$8\,056$}{$8\,560$}{$89$}
\cauTN{8}{Khoảng thời gian $5\,\text{phút}\,12\,\text{giây}$ bằng bao nhiêu giây?}{$300\,\text{giây}$}{$312\,\text{giây}$}{$322\,\text{giây}$}{$332\,\text{giây}$}
\cauTN{9}{Tìm $x$: $x - 47\,350 = 68\,490$:}{$x = 115\,840$}{$x = 21\,140$}{$x = 115\,740$}{$x = 116\,840$}
\cauTN{10}{Làm tròn số $846\,200$ đến hàng trăm nghìn ta được số:}{$800\,000$}{$900\,000$}{$850\,000$}{$840\,000$}
\cauTNhaiDong{11}{Tổng của hai số là $200$, hiệu của hai số là $60$. Số bé là:}{$60$}{$70$}{$80$}{$90$}
\cauTNhaiDong{12}{Một mảnh đất hình chữ nhật có chu vi $160\,\text{m}$, chiều dài hơn chiều rộng $24\,\text{m}$. Chiều rộng là:}{$28\,\text{m}$}{$52\,\text{m}$}{$30\,\text{m}$}{$35\,\text{m}$}

\vspace{0.4mm}
\noindent\textbf{\large PHẦN II. TỰ LUẬN (7,0 điểm)}

\noindent\textbf{Bài 1 (2,0 điểm). Đặt tính rồi tính:}
\vspace{0.4mm}

\noindent
\begin{tabularx}{\linewidth}{|X|X|X|X|}
\hline
\oDatTinh[4]{a)}{$563\,820 + 278\,695$} & 
\oDatTinh[4]{b)}{$842\,150 - 468\,735$} & 
\oDatTinh[4]{c)}{$16\,528 \times 5$} & 
\oDatTinh[4]{d)}{$59\,643 : 7$} \tabularnewline
\hline
\end{tabularx}

\newpage
% -------------------- TRANG 2: MẶT SAU --------------------
\noindent
\begin{tabularx}{\linewidth}{|X|c|}
\hline
\textbf{PHẦN TỰ LUẬN (TIẾP THEO) -- PHIẾU BÀI TẬP TUẦN 15} & \textbf{MÃ ĐỀ: 404} \\
\hline
\end{tabularx}

\vspace{0.8mm}

% BÀI 2: TÌM X (4 Ý)
\khungBaiTap{%
  \textbf{Bài 2 (1,5 điểm). Tìm $x$:}\par\vspace{0.6mm}
  \noindent
  \begin{minipage}[t]{0.485\linewidth}
    \textbf{a)} $x + 52\,470 = 98\,350$\par
    \dongke{3}
    \vspace{0.4mm}
    \textbf{c)} $98\,200 - x = 54\,630$\par
    \dongke{3}
  \end{minipage}\hfill
  \begin{minipage}[t]{0.485\linewidth}
    \textbf{b)} $x - 63\,280 = 79\,450$\par
    \dongke{3}
    \vspace{0.4mm}
    \textbf{d)} $x : 8 = 15\,348$\par
    \dongke{3}
  \end{minipage}
}

\vspace{0.8mm}

% BÀI 3: ĐỔI ĐƠN VỊ (4 Ý - 1 CỘT TỪ TRÊN XUỐNG)
\khungBaiTap{%
  \textbf{Bài 3 (1,0 điểm). Viết số thích hợp vào chỗ chấm:}\par\vspace{0.8mm}
  \noindent
  a) $7\,\text{tấn}\,650\,\text{kg} = \dotfill\,\text{kg}$\par\vspace{1.4mm}
  \noindent
  b) $9\,\text{m}^2\,35\,\text{dm}^2 = \dotfill\,\text{dm}^2$\par\vspace{1.4mm}
  \noindent
  c) $6\,\text{thế kỉ}\,45\,\text{năm} = \dotfill\,\text{năm}$\par\vspace{1.4mm}
  \noindent
  d) $3\,\text{phút}\,55\,\text{giây} = \dotfill\,\text{giây}$\par
}

\vspace{0.8mm}

% BÀI 4: BÀI TOÁN GIẢI (11 DÒNG KẺ)
\khungBaiTap{%
  \textbf{Bài 4 (1,5 điểm).} Một thửa đất hình thang có đáy lớn $35$ m, đáy bé $25$ m, chiều cao $20$ m. Người ta trồng rau trên thửa đất đó, cứ $5$ m$^2$ thu được $40\,000$ đồng. Tính số tiền thu được từ thửa đất.
  \par\vspace{0.3mm}%
  \noindent\textit{Bài giải:}
  \dongke{11}
}

\vspace{0.8mm}

% BÀI 5: TOÁN NÂNG CAO (8 DÒNG KẺ)
\khungBaiTap{%
  \textbf{Bài 5 (1,0 điểm).} Tính bằng cách thuận tiện nhất:\par\vspace{0.6mm}\centerline{$Tính bằng cách thuận tiện: D = 2{,}5 \times 4{,}8 + 2{,}5 \times 5{,}2$}
  \par\vspace{0.3mm}%
  \noindent\textit{Bài làm:}
  \dongke{8}
}

\newpage


% =============================================================================
% ========================= MÃ ĐỀ: 405 ========================================
% =============================================================================

% -------------------- TRANG 1: MẶT TRƯỚC --------------------
\noindent
\begin{tabularx}{\linewidth}{|X|c|}
\hline
\textbf{PHIẾU BÀI TẬP CUỐI TUẦN -- TOÁN LỚP 4 (TUẦN 15)} & \textbf{MÃ ĐỀ: 405} \\
\textit{Trung tâm Toán Thầy Lĩnh -- Năm học: 2026 - 2027} & \textit{Thời gian: 45 phút} \\
\hline
\end{tabularx}

\vspace{1mm}

\noindent
\begin{tabularx}{\linewidth}{|X|p{3.5cm}|}
\hline
\rule[-2mm]{0pt}{7.5mm}\textbf{Họ và tên:} \dotfill & \textbf{Lớp:} \makebox[2cm]{\dotfill} \\
\hline
\end{tabularx}

\vspace{1mm}

\noindent
\begin{tabularx}{\linewidth}{|p{3cm}|X|}
\hline
\centering\textbf{ĐIỂM SỐ} & \centering\textbf{LỜI NHẬN XÉT CỦA THẦY CÔ} \tabularnewline
\hline
\rule{0pt}{1.1cm} & \tabularnewline
\hline
\end{tabularx}

\vspace{1mm}

\noindent\textbf{BẢNG ĐIỀN ĐÁP ÁN TRẮC NGHIỆM (Mỗi câu đúng 0,25 điểm):}
\vspace{0.6mm}

\begin{tabularx}{\linewidth}{|c|*{12}{>{\centering\arraybackslash}X|}}
\hline
\textbf{Câu} & \textbf{1} & \textbf{2} & \textbf{3} & \textbf{4} & \textbf{5} & \textbf{6} & \textbf{7} & \textbf{8} & \textbf{9} & \textbf{10} & \textbf{11} & \textbf{12} \\
\hline
\textbf{Đáp án} & & & & & & & & & & & & \\
\hline
\end{tabularx}

\vspace{1.2mm}

\noindent\textbf{\large PHẦN I. TRẮC NGHIỆM KHÁCH QUAN (3,0 điểm)}

\cauTN{1}{Đường thẳng $a$ song song với đường thẳng $b$, ký hiệu là:}{$a \perp b$}{$a \parallel b$}{$a = b$}{$a + b$}
\cauTN{2}{Trong các hình sau, hình nào có $2$ cặp cạnh đối diện song song?}{Hình tam giác}{Hình thang vuông}{Hình chữ nhật}{Hình ngũ giác}
\cauTN{3}{Hai đường thẳng không cắt nhau trên một mặt phẳng gọi là hai đường thẳng:}{Vuông góc}{Song song}{Trùng nhau}{Chéo nhau}
\cauTN{4}{Hình thoi có mấy cặp cạnh song song?}{$1$ cặp}{$2$ cặp}{$3$ cặp}{$4$ cặp}
\cauTN{5}{Trong hình chữ nhật $ABCD$, cạnh $CD$ song song với cạnh:}{Cạnh $AB$}{Cạnh $BC$}{Cạnh $AD$}{Cạnh $BD$}
\cauTN{6}{Số thích hợp điền vào chỗ chấm $7\,\text{tấn}\,85\,\text{kg} = \dots\,\text{kg}$ là:}{$785$}{$7\,085$}{$7\,850$}{$78\,500$}
\cauTN{7}{Số thích hợp điền vào chỗ chấm $6\,\text{m}^2\,94\,\text{dm}^2 = \dots\,\text{dm}^2$ là:}{$694$}{$6\,094$}{$6\,940$}{$68$}
\cauTN{8}{Khoảng thời gian $6\,\text{phút}\,25\,\text{giây}$ bằng bao nhiêu giây?}{$360\,\text{giây}$}{$375\,\text{giây}$}{$385\,\text{giây}$}{$395\,\text{giây}$}
\cauTN{9}{Tìm $x$: $x + 46\,720 = 91\,580$:}{$x = 44\,860$}{$x = 44\,960$}{$x = 138\,300$}{$x = 43\,860$}
\cauTN{10}{Làm tròn số $935\,600$ đến hàng trăm nghìn ta được số:}{$900\,000$}{$1\,000\,000$}{$930\,000$}{$940\,000$}
\cauTNhaiDong{11}{Tổng của hai số là $220$, hiệu của hai số là $40$. Số lớn là:}{$120$}{$130$}{$140$}{$150$}
\cauTNhaiDong{12}{Một mảnh đất hình chữ nhật có chu vi $180\,\text{m}$, chiều dài hơn chiều rộng $26\,\text{m}$. Chiều dài là:}{$58\,\text{m}$}{$32\,\text{m}$}{$60\,\text{m}$}{$45\,\text{m}$}

\vspace{0.4mm}
\noindent\textbf{\large PHẦN II. TỰ LUẬN (7,0 điểm)}

\noindent\textbf{Bài 1 (2,0 điểm). Đặt tính rồi tính:}
\vspace{0.4mm}

\noindent
\begin{tabularx}{\linewidth}{|X|X|X|X|}
\hline
\oDatTinh[4]{a)}{$478\,530 + 386\,794$} & 
\oDatTinh[4]{b)}{$924\,150 - 468\,725$} & 
\oDatTinh[4]{c)}{$15\,486 \times 6$} & 
\oDatTinh[4]{d)}{$68\,544 : 8$} \tabularnewline
\hline
\end{tabularx}

\newpage
% -------------------- TRANG 2: MẶT SAU --------------------
\noindent
\begin{tabularx}{\linewidth}{|X|c|}
\hline
\textbf{PHẦN TỰ LUẬN (TIẾP THEO) -- PHIẾU BÀI TẬP TUẦN 15} & \textbf{MÃ ĐỀ: 405} \\
\hline
\end{tabularx}

\vspace{0.8mm}

% BÀI 2: TÌM X (4 Ý)
\khungBaiTap{%
  \textbf{Bài 2 (1,5 điểm). Tìm $x$:}\par\vspace{0.6mm}
  \noindent
  \begin{minipage}[t]{0.485\linewidth}
    \textbf{a)} $x + 53\,480 = 97\,250$\par
    \dongke{3}
    \vspace{0.4mm}
    \textbf{c)} $96\,500 - x = 52\,380$\par
    \dongke{3}
  \end{minipage}\hfill
  \begin{minipage}[t]{0.485\linewidth}
    \textbf{b)} $x - 48\,650 = 85\,720$\par
    \dongke{3}
    \vspace{0.4mm}
    \textbf{d)} $x : 7 = 18\,345$\par
    \dongke{3}
  \end{minipage}
}

\vspace{0.8mm}

% BÀI 3: ĐỔI ĐƠN VỊ (4 Ý - 1 CỘT TỪ TRÊN XUỐNG)
\khungBaiTap{%
  \textbf{Bài 3 (1,0 điểm). Viết số thích hợp vào chỗ chấm:}\par\vspace{0.8mm}
  \noindent
  a) $8\,\text{tấn}\,540\,\text{kg} = \dotfill\,\text{kg}$\par\vspace{1.4mm}
  \noindent
  b) $7\,\text{m}^2\,62\,\text{dm}^2 = \dotfill\,\text{dm}^2$\par\vspace{1.4mm}
  \noindent
  c) $5\,\text{thế kỉ}\,80\,\text{năm} = \dotfill\,\text{năm}$\par\vspace{1.4mm}
  \noindent
  d) $4\,\text{phút}\,30\,\text{giây} = \dotfill\,\text{giây}$\par
}

\vspace{0.8mm}

% BÀI 4: BÀI TOÁN GIẢI (11 DÒNG KẺ)
\khungBaiTap{%
  \textbf{Bài 4 (1,5 điểm).} Một sân hình chữ nhật có chiều dài $40$ m, chiều rộng $25$ m. Người ta làm một bồn hoa hình tròn đường kính $10$ m ở giữa sân (lấy $\pi = 3{,}14$). Phần còn lại lát gạch với giá $120\,000$ đồng mỗi mét vuông. Tính số tiền lát gạch.
  \par\vspace{0.3mm}%
  \noindent\textit{Bài giải:}
  \dongke{11}
}

\vspace{0.8mm}

% BÀI 5: TOÁN NÂNG CAO (8 DÒNG KẺ)
\khungBaiTap{%
  \textbf{Bài 5 (1,0 điểm).} Tính bằng cách thuận tiện nhất:\par\vspace{0.6mm}\centerline{$Tính bằng cách thuận tiện: E = 0{,}5 + 0{,}5 + 0{,}5 + 0{,}5 + 0{,}5 + 0{,}5$}
  \par\vspace{0.3mm}%
  \noindent\textit{Bài làm:}
  \dongke{8}
}

\newpage


% =============================================================================
% ========================= MÃ ĐỀ: 406 ========================================
% =============================================================================

% -------------------- TRANG 1: MẶT TRƯỚC --------------------
\noindent
\begin{tabularx}{\linewidth}{|X|c|}
\hline
\textbf{PHIẾU BÀI TẬP CUỐI TUẦN -- TOÁN LỚP 4 (TUẦN 15)} & \textbf{MÃ ĐỀ: 406} \\
\textit{Trung tâm Toán Thầy Lĩnh -- Năm học: 2026 - 2027} & \textit{Thời gian: 45 phút} \\
\hline
\end{tabularx}

\vspace{1mm}

\noindent
\begin{tabularx}{\linewidth}{|X|p{3.5cm}|}
\hline
\rule[-2mm]{0pt}{7.5mm}\textbf{Họ và tên:} \dotfill & \textbf{Lớp:} \makebox[2cm]{\dotfill} \\
\hline
\end{tabularx}

\vspace{1mm}

\noindent
\begin{tabularx}{\linewidth}{|p{3cm}|X|}
\hline
\centering\textbf{ĐIỂM SỐ} & \centering\textbf{LỜI NHẬN XÉT CỦA THẦY CÔ} \tabularnewline
\hline
\rule{0pt}{1.1cm} & \tabularnewline
\hline
\end{tabularx}

\vspace{1mm}

\noindent\textbf{BẢNG ĐIỀN ĐÁP ÁN TRẮC NGHIỆM (Mỗi câu đúng 0,25 điểm):}
\vspace{0.6mm}

\begin{tabularx}{\linewidth}{|c|*{12}{>{\centering\arraybackslash}X|}}
\hline
\textbf{Câu} & \textbf{1} & \textbf{2} & \textbf{3} & \textbf{4} & \textbf{5} & \textbf{6} & \textbf{7} & \textbf{8} & \textbf{9} & \textbf{10} & \textbf{11} & \textbf{12} \\
\hline
\textbf{Đáp án} & & & & & & & & & & & & \\
\hline
\end{tabularx}

\vspace{1.2mm}

\noindent\textbf{\large PHẦN I. TRẮC NGHIỆM KHÁCH QUAN (3,0 điểm)}

\cauTN{1}{Hai đường thẳng $a$ và $b$ cùng song song với đường thẳng $c$ thì:}{$a$ vuông góc với $b$}{$a$ song song với $b$}{$a$ cắt $b$}{Không kết luận được}
\cauTN{2}{Trong hình vuông $MNPQ$, cặp cạnh nào sau đây song song với nhau?}{$MN$ và $NP$}{$MN$ và $PQ$}{$NP$ và $PQ$}{$MQ$ và $MN$}
\cauTN{3}{Hình tứ giác có hai cặp cạnh đối diện song song và bằng nhau là:}{Hình chữ nhật (hoặc hình bình hành)}{Hình tam giác}{Hình thang}{Hình tròn}
\cauTN{4}{Độ dài các đoạn thẳng vuông góc nối giữa hai đường thẳng song song thì:}{Bằng nhau}{Không bằng nhau}{Tăng dần}{Giảm dần}
\cauTN{5}{Cạnh đối diện của hình chữ nhật dài $15\,\text{cm}$ thì cạnh kia dài:}{$10\,\text{cm}$}{$15\,\text{cm}$}{$20\,\text{cm}$}{$30\,\text{cm}$}
\cauTN{6}{Số thích hợp điền vào chỗ chấm $9\,\text{tấn}\,45\,\text{kg} = \dots\,\text{kg}$ là:}{$945$}{$9\,045$}{$9\,450$}{$94\,500$}
\cauTN{7}{Số thích hợp điền vào chỗ chấm $8\,\text{m}^2\,85\,\text{dm}^2 = \dots\,\text{dm}^2$ là:}{$885$}{$8\,085$}{$8\,850$}{$89$}
\cauTN{8}{Khoảng thời gian $5\,\text{phút}\,40\,\text{giây}$ bằng bao nhiêu giây?}{$320\,\text{giây}$}{$340\,\text{giây}$}{$350\,\text{giây}$}{$360\,\text{giây}$}
\cauTN{9}{Tìm $x$: $x - 56\,380 = 78\,450$:}{$x = 134\,830$}{$x = 22\,070$}{$x = 134\,730$}{$x = 135\,830$}
\cauTN{10}{Làm tròn số $854\,700$ đến hàng trăm nghìn ta được số:}{$800\,000$}{$900\,000$}{$850\,000$}{$860\,000$}
\cauTNhaiDong{11}{Tổng của hai số là $240$, hiệu của hai số là $60$. Số bé là:}{$80$}{$90$}{$100$}{$110$}
\cauTNhaiDong{12}{Một mảnh đất hình chữ nhật có chu vi $220\,\text{m}$, chiều dài hơn chiều rộng $30\,\text{m}$. Chiều rộng là:}{$40\,\text{m}$}{$70\,\text{m}$}{$35\,\text{m}$}{$45\,\text{m}$}

\vspace{0.4mm}
\noindent\textbf{\large PHẦN II. TỰ LUẬN (7,0 điểm)}

\noindent\textbf{Bài 1 (2,0 điểm). Đặt tính rồi tính:}
\vspace{0.4mm}

\noindent
\begin{tabularx}{\linewidth}{|X|X|X|X|}
\hline
\oDatTinh[4]{a)}{$642\,830 + 279\,585$} & 
\oDatTinh[4]{b)}{$863\,400 - 487\,562$} & 
\oDatTinh[4]{c)}{$17\,346 \times 5$} & 
\oDatTinh[4]{d)}{$73\,656 : 9$} \tabularnewline
\hline
\end{tabularx}

\newpage
% -------------------- TRANG 2: MẶT SAU --------------------
\noindent
\begin{tabularx}{\linewidth}{|X|c|}
\hline
\textbf{PHẦN TỰ LUẬN (TIẾP THEO) -- PHIẾU BÀI TẬP TUẦN 15} & \textbf{MÃ ĐỀ: 406} \\
\hline
\end{tabularx}

\vspace{0.8mm}

% BÀI 2: TÌM X (4 Ý)
\khungBaiTap{%
  \textbf{Bài 2 (1,5 điểm). Tìm $x$:}\par\vspace{0.6mm}
  \noindent
  \begin{minipage}[t]{0.485\linewidth}
    \textbf{a)} $x + 47\,850 = 96\,320$\par
    \dongke{3}
    \vspace{0.4mm}
    \textbf{c)} $97\,500 - x = 53\,480$\par
    \dongke{3}
  \end{minipage}\hfill
  \begin{minipage}[t]{0.485\linewidth}
    \textbf{b)} $x - 58\,420 = 86\,390$\par
    \dongke{3}
    \vspace{0.4mm}
    \textbf{d)} $x : 8 = 19\,245$\par
    \dongke{3}
  \end{minipage}
}

\vspace{0.8mm}

% BÀI 3: ĐỔI ĐƠN VỊ (4 Ý - 1 CỘT TỪ TRÊN XUỐNG)
\khungBaiTap{%
  \textbf{Bài 3 (1,0 điểm). Viết số thích hợp vào chỗ chấm:}\par\vspace{0.8mm}
  \noindent
  a) $9\,\text{tấn}\,360\,\text{kg} = \dotfill\,\text{kg}$\par\vspace{1.4mm}
  \noindent
  b) $8\,\text{m}^2\,47\,\text{dm}^2 = \dotfill\,\text{dm}^2$\par\vspace{1.4mm}
  \noindent
  c) $7\,\text{thế kỉ}\,60\,\text{năm} = \dotfill\,\text{năm}$\par\vspace{1.4mm}
  \noindent
  d) $6\,\text{phút}\,10\,\text{giây} = \dotfill\,\text{giây}$\par
}

\vspace{0.8mm}

% BÀI 4: BÀI TOÁN GIẢI (11 DÒNG KẺ)
\khungBaiTap{%
  \textbf{Bài 4 (1,5 điểm).} Một mảnh đất hình thang có đáy lớn $28$ m, đáy bé $12$ m, chiều cao $15$ m. Người ta chia mảnh đất thành hai phần: phần hình tam giác có đáy$12$ m (bằng đáy bé) và chiều cao $15$ m để làm sân, phần còn lại trồng hoa với giá $30\,000$ đồng mỗi mét vuông. Tính số tiền trồng hoa.
  \par\vspace{0.3mm}%
  \noindent\textit{Bài giải:}
  \dongke{11}
}

\vspace{0.8mm}

% BÀI 5: TOÁN NÂNG CAO (8 DÒNG KẺ)
\khungBaiTap{%
  \textbf{Bài 5 (1,0 điểm).} Tính bằng cách thuận tiện nhất:\par\vspace{0.6mm}\centerline{$Tính bằng cách thuận tiện: F = 0{,}12 \times 5 + 0{,}12 \times 95$}
  \par\vspace{0.3mm}%
  \noindent\textit{Bài làm:}
  \dongke{8}
}

\newpage


% =============================================================================
% ========================= MÃ ĐỀ: 407 ========================================
% =============================================================================

% -------------------- TRANG 1: MẶT TRƯỚC --------------------
\noindent
\begin{tabularx}{\linewidth}{|X|c|}
\hline
\textbf{PHIẾU BÀI TẬP CUỐI TUẦN -- TOÁN LỚP 4 (TUẦN 15)} & \textbf{MÃ ĐỀ: 407} \\
\textit{Trung tâm Toán Thầy Lĩnh -- Năm học: 2026 - 2027} & \textit{Thời gian: 45 phút} \\
\hline
\end{tabularx}

\vspace{1mm}

\noindent
\begin{tabularx}{\linewidth}{|X|p{3.5cm}|}
\hline
\rule[-2mm]{0pt}{7.5mm}\textbf{Họ và tên:} \dotfill & \textbf{Lớp:} \makebox[2cm]{\dotfill} \\
\hline
\end{tabularx}

\vspace{1mm}

\noindent
\begin{tabularx}{\linewidth}{|p{3cm}|X|}
\hline
\centering\textbf{ĐIỂM SỐ} & \centering\textbf{LỜI NHẬN XÉT CỦA THẦY CÔ} \tabularnewline
\hline
\rule{0pt}{1.1cm} & \tabularnewline
\hline
\end{tabularx}

\vspace{1mm}

\noindent\textbf{BẢNG ĐIỀN ĐÁP ÁN TRẮC NGHIỆM (Mỗi câu đúng 0,25 điểm):}
\vspace{0.6mm}

\begin{tabularx}{\linewidth}{|c|*{12}{>{\centering\arraybackslash}X|}}
\hline
\textbf{Câu} & \textbf{1} & \textbf{2} & \textbf{3} & \textbf{4} & \textbf{5} & \textbf{6} & \textbf{7} & \textbf{8} & \textbf{9} & \textbf{10} & \textbf{11} & \textbf{12} \\
\hline
\textbf{Đáp án} & & & & & & & & & & & & \\
\hline
\end{tabularx}

\vspace{1.2mm}

\noindent\textbf{\large PHẦN I. TRẮC NGHIỆM KHÁCH QUAN (3,0 điểm)}

\cauTN{1}{Đặc điểm nào sau đây KHÔNG đúng với hai đường thẳng song song?}{Không bao giờ cắt nhau}{Cắt nhau tạo thành 4 góc vuông}{Khoảng cách giữa chúng không đổi}{Cùng vuông góc với một đường thẳng thứ ba}
\cauTN{2}{Trong hình bình hành $ABCD$, cạnh $AB$ song song với cạnh nào?}{Cạnh $BC$}{Cạnh $CD$}{Cạnh $AD$}{Cạnh $AC$}
\cauTN{3}{Một hình vuông có cạnh dài $8\,\text{cm}$, chu vi hình vuông đó là:}{$16\,\text{cm}$}{$24\,\text{cm}$}{$32\,\text{cm}$}{$64\,\text{cm}^2$}
\cauTN{4}{Hai cạnh đối diện của bảng lớp học là hình ảnh của:}{Hai đường thẳng vuông góc}{Hai đường thẳng song song}{Hai đường thẳng cắt nhau}{Hai góc bẹt}
\cauTN{5}{Hai mép đối diện của thước kẻ thẳng là hình ảnh của:}{Hai đường thẳng song song}{Hai đường thẳng vuông góc}{Hai đường cong}{Hai điểm}
\cauTN{6}{Số thích hợp điền vào chỗ chấm $6\,\text{tấn}\,95\,\text{kg} = \dots\,\text{kg}$ là:}{$695$}{$6\,095$}{$6\,950$}{$69\,500$}
\cauTN{7}{Số thích hợp điền vào chỗ chấm $9\,\text{m}^2\,18\,\text{dm}^2 = \dots\,\text{dm}^2$ là:}{$918$}{$9\,018$}{$9\,180$}{$92$}
\cauTN{8}{Khoảng thời gian $6\,\text{phút}\,45\,\text{giây}$ bằng bao nhiêu giây?}{$385\,\text{giây}$}{$395\,\text{giây}$}{$405\,\text{giây}$}{$415\,\text{giây}$}
\cauTN{9}{Tìm $x$: $x + 63\,480 = 98\,250$:}{$x = 34\,770$}{$x = 34\,870$}{$x = 161\,730$}{$x = 33\,770$}
\cauTN{10}{Làm tròn số $972\,500$ đến hàng trăm nghìn ta được số:}{$900\,000$}{$970\,000$}{$1\,000\,000$}{$980\,000$}
\cauTNhaiDong{11}{Tổng của hai số là $260$, hiệu của hai số là $80$. Số lớn là:}{$160$}{$170$}{$180$}{$190$}
\cauTNhaiDong{12}{Một mảnh đất hình chữ nhật có chu vi $240\,\text{m}$, chiều dài hơn chiều rộng $32\,\text{m}$. Chiều rộng là:}{$44\,\text{m}$}{$76\,\text{m}$}{$45\,\text{m}$}{$50\,\text{m}$}

\vspace{0.4mm}
\noindent\textbf{\large PHẦN II. TỰ LUẬN (7,0 điểm)}

\noindent\textbf{Bài 1 (2,0 điểm). Đặt tính rồi tính:}
\vspace{0.4mm}

\noindent
\begin{tabularx}{\linewidth}{|X|X|X|X|}
\hline
\oDatTinh[4]{a)}{$725\,640 + 198\,785$} & 
\oDatTinh[4]{b)}{$942\,300 - 478\,654$} & 
\oDatTinh[4]{c)}{$18\,456 \times 5$} & 
\oDatTinh[4]{d)}{$84\,672 : 9$} \tabularnewline
\hline
\end{tabularx}

\newpage
% -------------------- TRANG 2: MẶT SAU --------------------
\noindent
\begin{tabularx}{\linewidth}{|X|c|}
\hline
\textbf{PHẦN TỰ LUẬN (TIẾP THEO) -- PHIẾU BÀI TẬP TUẦN 15} & \textbf{MÃ ĐỀ: 407} \\
\hline
\end{tabularx}

\vspace{0.8mm}

% BÀI 2: TÌM X (4 Ý)
\khungBaiTap{%
  \textbf{Bài 2 (1,5 điểm). Tìm $x$:}\par\vspace{0.6mm}
  \noindent
  \begin{minipage}[t]{0.485\linewidth}
    \textbf{a)} $x + 58\,420 = 97\,350$\par
    \dongke{3}
    \vspace{0.4mm}
    \textbf{c)} $98\,600 - x = 56\,370$\par
    \dongke{3}
  \end{minipage}\hfill
  \begin{minipage}[t]{0.485\linewidth}
    \textbf{b)} $x - 64\,390 = 87\,560$\par
    \dongke{3}
    \vspace{0.4mm}
    \textbf{d)} $x : 9 = 17\,348$\par
    \dongke{3}
  \end{minipage}
}

\vspace{0.8mm}

% BÀI 3: ĐỔI ĐƠN VỊ (4 Ý - 1 CỘT TỪ TRÊN XUỐNG)
\khungBaiTap{%
  \textbf{Bài 3 (1,0 điểm). Viết số thích hợp vào chỗ chấm:}\par\vspace{0.8mm}
  \noindent
  a) $9\,\text{tấn}\,450\,\text{kg} = \dotfill\,\text{kg}$\par\vspace{1.4mm}
  \noindent
  b) $7\,\text{m}^2\,86\,\text{dm}^2 = \dotfill\,\text{dm}^2$\par\vspace{1.4mm}
  \noindent
  c) $8\,\text{thế kỉ}\,50\,\text{năm} = \dotfill\,\text{năm}$\par\vspace{1.4mm}
  \noindent
  d) $5\,\text{phút}\,55\,\text{giây} = \dotfill\,\text{giây}$\par
}

\vspace{0.8mm}

% BÀI 4: BÀI TOÁN GIẢI (11 DÒNG KẺ)
\khungBaiTap{%
  \textbf{Bài 4 (1,5 điểm).} Một khu đất hình chữ nhật có chiều dài $50$ m, chiều rộng $30$ m. Người ta dùng $\frac{2}{5}$ diện tích để xây nhà, phần còn lại chia đều thành các ô đất hình tam giác có đáy $4$ m, chiều cao $5$ m để trồng cây. Mỗi ô đất trồng cây tốn $200\,000$ đồng. Tính tổng số tiền trồng cây.
  \par\vspace{0.3mm}%
  \noindent\textit{Bài giải:}
  \dongke{11}
}

\vspace{0.8mm}

% BÀI 5: TOÁN NÂNG CAO (8 DÒNG KẺ)
\khungBaiTap{%
  \textbf{Bài 5 (1,0 điểm).} Tính bằng cách thuận tiện nhất:\par\vspace{0.6mm}\centerline{$Tính bằng cách thuận tiện: G = 0{,}25 \times 1{,}2 + 0{,}25 \times 8{,}8$}
  \par\vspace{0.3mm}%
  \noindent\textit{Bài làm:}
  \dongke{8}
}

\newpage
\end{document}
